\section{Finite computations}
\label{supp:evidence}

The computations below are exact (integer, rational or algebraic arithmetic, or a SAT solver) and
finite. None of them is used in a proof in the main paper; each illustrates a law on a bounded
range, and a bounded run is evidence only about that range. The parameters of every run are stored
with its output; \Cref{supp:replay} gives the commands.

{\small
\begin{longtable}{@{}l>{\raggedright\arraybackslash}p{5.6cm}>{\raggedright\arraybackslash}p{7.6cm}@{}}
\caption{Finite computations, by law.}\label{tab:evidence}\\
\toprule
\textbf{Run} & \textbf{Computation} & \textbf{Outcome} \\
\midrule
\endfirsthead
\toprule
\textbf{Run} & \textbf{Computation} & \textbf{Outcome} \\
\midrule
\endhead
\bottomrule
\endfoot
G1-L1 & Accelerated Collatz orbits of all $500000$ odd $n < 10^6$, with the ledger identity of
\Cref{thm:affine-ledger} checked at $2241405$ checkpoints. & Every start other than $1$ descends;
the latest first descent is at step $111$, for $n = 626331$. \\
G1-L2 & All odd starts below $10^7$; the $2^L - 1$ formula of \Cref{rem:collatz-paper} for
$L = 2, \ldots, 80$. & Latest first descent at step $155$, for $n = 8088063$. The formula holds
for every $L$ tested. A stronger guess, that the first descent of $2^L - 1$ occurs exactly at step
$L - 1$, failed in all $79$ cases (for $L = 2$, $n = 3$ descends at step $2$); the main paper
claims only ``at least $L - 1$''. \\
G1-L3 & A finite analogue of the ghost: the accelerated map on the $32768$ odd residues modulo
$2^{16}$. & Apart from the fixed residue $1$, every residue descends except $14563$ and $21845$, which
never descend and reach the ghost of this finite model, the odd solution $21845 = (2^{16} - 1)/3$
of $3r + 1 \equiv 0 \pmod{2^{16}}$. In this finite model both hypotheses of \Cref{thm:G1b} hold;
this says nothing about the integers. \\
G1-L4 & Aliquot sequences from $100000$ starts with a cap of $10^{12}$. & $82728$ reach $0$, $2085$
enter one of $22$ cycles, and $15187$ exceed the cap. \\
G2-L1 & Goodstein sequences from $1, \ldots, 30$ for $4$ steps, checking that rewriting preserves
the ordinal and subtraction lowers it. & All $116$ steps satisfy both identities; the largest value
reached has $120605$ bits. Starts $1$ and $2$ terminate within the budget. \\
G3-L1 & $1000$ binary-counter increments from $44142$; $1000$ dynamic-array insertions. & Amortized
cost exactly $2$ per increment; at most $3$ per insertion, with equality at each of the $8$
resizes. \\
G4-L1 & Greedy unit-fraction expansions of all $124750$ fractions $p/q$ with $0 < p < q \le 500$. &
Every expansion terminates with strictly decreasing numerators. \\
G5-L1 & The base-$2$ reverse-and-add orbit of $22$ for $84$ steps. & Every phase transition of
\Cref{thm:four-phase} holds; the largest word has $48$ bits. \\
G6-L1 & $10^7$ steps of Recam\'an's sequence. & The backward move is legal at $4999986$ steps;
$852655$ is still missing. \\
G6-L2 & The two designed variants of \Cref{sec:G6} for $10^6$ steps. & The variant with jumps $2n$
visits $736749$ even values and no odd value; the unit-jump variant visits exactly
$\{0, \ldots, n\}$ by step $n$. \\
G7-L1 & Bounded-horizon minimax for the angel game on small tori, $8$ parameter settings. & On the
$4 \times 4$ torus with $b = 2$, power $1$ is trapped and power $2$ survives to the horizon; the
$3 \times 3$ settings cannot distinguish powers $1$ and $2$, and $b = 1$ on $4 \times 4$ was
inconclusive. A finite torus says nothing about the infinite board. \\
G8-L1 & $500$ stabilizations of random sandpiles on square grids of sides $10$ to $50$, $100$
random legal orders each. & Every order gives the same final configuration and odometer. \\
G8-L2 & $1000$ random legal orders in \Cref{ex:G8-b}. & Every run ends at $N$; exactly three
counter vectors occur. \\
G9-L1 & Rule~184 on a ring of $2000$ cells at densities $\tfrac{3}{10}, \tfrac25, \tfrac12,
\tfrac35, \tfrac{7}{10}$. & Each run reaches the predicted regime; the recurrence certificate
holds on $128$ consecutive windows, with period $1$ and displacement $+1$ below density $\tfrac12$
and $-1$ above. \\
G10-L1 & Integer Ducci tuples of lengths $3$ to $16$. & Lengths $4$, $8$, $16$ reach zero (at steps
$4$, $18$, $98$ in the recorded runs), and $(I + S)^k = 0$ over $\mathbb{F}_2$ exactly for these
lengths; the tuples of the other lengths enter nonzero cycles. \\
G11-L1 & SAT encodings of periodic tilings on tori. & A one-colour tile set tiles the
$1 \times 1$ torus; the $11$-tile set of \citet{JeandelRao2021} tiles no $p_1 \times p_2$ torus with
$p_1, p_2 \le 4$. \\
G12-L1 & An $826$-vertex unit-distance graph from the Polymath16 project, in exact algebraic
arithmetic. & All $4273$ edges have unit length and $2000$ sampled non-edges do not; the graph is
not $4$-colourable (SAT, $3304$ variables, $22874$ clauses). No independently checkable proof of
unsatisfiability is available for this graph. \\
G13-L1 & Apollonian packings from $(-1, 2, 2, 3)$ and $(-6, 11, 14, 15)$ up to curvature
$300000$. & Descartes' relation holds for every quadruple, and all curvatures lie in the expected
residue classes modulo $24$. The second packing contains no curvature of the form $2n^2$, $3n^2$ or
$6n^2$ in range, as the reciprocity obstruction predicts. \\
\end{longtable}
}
